\section{Introduction}
Fairness testing is usually presented as an audit of one fixed system. Software, however, evolves: a previously certified program $P_0$ is replaced by a patched program $P_1$, and the auditor must decide whether the change has introduced a forbidden discrepancy. When execution is stochastic, evaluating the same input twice may not give the same result. The central question becomes not only \emph{which comparisons} remain necessary but also \emph{which old/new executions can share randomness}. The work in \cite{BuselliEtAl2024} uses ``unfair regression'' for sensitive-group imbalance in negative prediction flips after an update; here a \emph{fairness regression} is instead a new violation on a declared counterfactual edge.

We formulate a query-complexity problem on a counterfactual graph $G=(V,E)$ whose vertices are valid program inputs and whose edges encode fairness comparisons. A sound impact set or sound effect envelope can eliminate edges without reexecution. For the residual graph, a tester chooses among old/new paired calls, four-way edge replay, full-suite replay, and restricted shared-randomness groups. These are different information resources, and a replay \emph{round} is not interchangeable with a concrete program execution.

The evidence-reuse workflow is summarized in \Cref{fig:pipeline}. The motivation draws on fairness testing \cite{GalhotraBrunMeliou2017,UdeshiAroraChattopadhyay2018,AlbarghouthiEtAl2017}, regression selection and nondeterministic testing \cite{RothermelHarrold1996,YooHarman2012,YuanLinShi2026}, and common-random-number coupling \cite{GlassermanYao1992,BartheEtAl2017}. The central theoretical question concerns the \emph{interaction} among these ingredients: how topology, software change, statistical uncertainty, and legal replay actions jointly determine testing cost. We do not claim that coupling, fractional covering, or sequential likelihood-ratio testing is itself new.

\paragraph{Main message.} Under tightly specified random-tape laws, stronger replay can change the dependence on a small fairness regression from $d^{-2}$ to $d^{-1}$. The improvement is not universal: positive tolerance, heterogeneous benign patch effects, and unavailable controls can remove it. Conversely, an appropriate control topology can restore the first-order information channel. Theorems and lower bounds therefore must explicitly name the observation law, comparator family, null class, and cost unit.

\paragraph{Principal results.} First, the Gaussian delta oracle induces an uncertainty-weighted graph allocation whose static upper and adaptive null-side lower bounds match up to confidence constants and a logarithmic familywise factor (\Cref{thm:upper,thm:lower}). Second, raw common-uniform replay yields a linear $d^{-1}$ information channel for sparse Bernoulli changes, while fixed positive tolerance restores quadratic $d^{-2}$ scaling (\Cref{thm:bernoulli-separation,thm:threshold-replay-phase}). Third, partial replay separates edge-witness and active-vertex cover geometries; in the declared two-class family, matched-background controls recover the linear channel, whereas boundary-only channels retain quadratic cost and exhibit a fixed-allocation familywise penalty removed from expected null cost by sequential stopping (\Cref{thm:control-replay,prop:boundary-familywise-fixed}). Certificate-aware localization and the executable historical defect connect these results to evolving software. Detailed proofs, counterexamples, and the experimental protocol follow.

\begin{figure}[!htbp]
\centering
\begin{tikzpicture}[
  >=Latex,
  every node/.style={font=\footnotesize},
  stage/.style={draw=linegray, line width=0.65pt, rounded corners=1.5pt,
    fill=white, align=left, inner xsep=6pt, inner ysep=6pt,
    text width=3.18cm, minimum height=1.42cm},
  focal/.style={stage, draw=focalteal, line width=0.9pt, fill=focalteal!8},
  model/.style={draw=linegray, line width=0.55pt, rounded corners=1.5pt,
    fill=boxedgray, align=left, inner xsep=6pt, inner ysep=5pt,
    text width=4.48cm, minimum height=1.28cm},
  flow/.style={-{Latex[length=2.0mm]}, draw=black!68, line width=0.72pt},
  choice/.style={-{Latex[length=1.8mm]}, draw=focalteal, line width=0.72pt},
  badge/.style={circle, draw=none, fill=focalteal, text=white,
    font=\footnotesize\bfseries, inner sep=1.7pt}
]
% Built at the final 16.0 cm width of the single-column review manuscript.
\node[stage] (archive) at (0,0) {\hspace{12pt}\textbf{Reuse certified evidence}\\[-1pt]
  $P_0$, patch $P_1$, fairness graph, and sound change information};
\node[focal] (reduce) at (4.12,0) {\hspace{12pt}\textbf{Remove safe obligations}\\[-1pt]
  Sound envelopes prune edges and produce residual graph $H$};
\node[focal] (match) at (8.24,0) {\hspace{12pt}\textbf{Match the \mbox{replay law}}\\[-1pt]
  Declare shared tapes, retained outputs, null class, and cost unit};
\node[stage] (report) at (12.36,0) {\hspace{12pt}\textbf{\mbox{Certify and report}}\\[-1pt]
  Apply only a scope-matched theorem; record decision and execution cost};
\node[badge, anchor=center] at ([xshift=7pt,yshift=-7pt]archive.north west) {1};
\node[badge, anchor=center] at ([xshift=7pt,yshift=-7pt]reduce.north west) {2};
\node[badge, anchor=center] at ([xshift=7pt,yshift=-7pt]match.north west) {3};
\node[badge, anchor=center] at ([xshift=7pt,yshift=-7pt]report.north west) {4};

\draw[flow] (archive.east) -- (reduce.west);
\draw[flow] (reduce.east) -- (match.west);
\draw[flow] (match.east) -- (report.west);

\node[model] (gaussian) at (0.20,-2.38) {\textbf{Paired Gaussian deltas}\\[-1pt]
  Vertex reuse yields the graph allocation $\Psi(H,\nu)$.};
\node[model] (uniform) at (5.48,-2.38) {\textbf{Raw synchronized replay}\\[-1pt]
  Sparse witnesses give $d^{-1}$; positive tolerance can restore $d^{-2}$.};
\node[model] (groups) at (10.76,-2.38) {\textbf{Legal partial-replay groups}\\[-1pt]
  Witness and matched-control covers encode which executions may couple.};

\coordinate (fork) at (8.24,-1.27);
\draw[choice] (match.south) -- (fork);
\draw[choice] (fork) -| (gaussian.north);
\draw[choice] (fork) -| (uniform.north);
\draw[choice] (fork) -| (groups.north);

\node[anchor=north, align=center, text=black!62, font=\footnotesize,
  text width=15.2cm] at (6.18,-4.05)
  {The lower boxes are distinct observation models---equal execution counts do not make them interchangeable.};
\end{tikzpicture}
\caption{Evidence reuse precedes statistical testing. Certified information first removes provably safe comparisons; each residual obligation is then matched to a declared replay law and to a theorem proved for that law. The workflow prevents a bound for one observation model or null family from being transferred silently to another.}
\Description{Four-stage workflow from certified old-version evidence and a software patch, through sound pruning and replay-law selection, to a scope-matched fairness decision. Three branches distinguish paired Gaussian deltas, synchronized Bernoulli replay, and legal partial-replay groups.}
\label{fig:pipeline}
\end{figure}


\section{Model and oracle hierarchy}
For each version $j\in\{0,1\}$ let $P_j:\mathcal X\times\mathcal R\to\mathcal Y$ be executable under random tape $R$. Write $\mu_{j,v}=\E[P_j(v;R)]$. A fairness edge $e=(u,v)$ requires $|\mu_{j,u}-\mu_{j,v}|\le\tau_e$ for a tolerance $\tau_e\ge0$. The old version is certified; $P_1$ must be tested. Let $\theta_v=\mu_{1,v}-\mu_{0,v}$ and $g_{j,e}=\mu_{j,u}-\mu_{j,v}$. Then
\begin{equation}\label{eq:change-identity}
g_{1,e}=g_{0,e}+\theta_u-\theta_v.
\end{equation}
A distributionally sound impact set $C_\Delta\subseteq V$ guarantees unchanged output laws outside it; a stronger replay-sound set guarantees pathwise old/new equality under a shared tape. All newly violated edges must meet $C_\Delta$. More generally, an upper bound $b_e$ on $|\theta_u-\theta_v|$ certifies edge $e$ without execution when $|g_{0,e}|+b_e\le\tau_e$.

\paragraph{Calls, rounds, and cost.} A single-version call executes one program version once. A vertex-paired execution runs $P_0(v;R)$ and $P_1(v;R)$ at one shared tape, costing two concrete executions. Its \emph{raw-pair} observation retains both outputs; the Gaussian \emph{delta oracle} instead reveals only their difference. Four-way, suite, and group replay retain all ordered raw outputs generated on their common tape. Fresh rounds use independent tapes. Four-way replay costs four executions; full-suite replay costs $2n_+$ per round; and a legal group $S$ costs $2|S|$ unless a different charge $c(S)$ is declared. Equal execution counts do not make these observation sigma-algebras interchangeable.

\paragraph{Core versus full fairness.} A Gaussian subproblem tests $H_0:\theta_u=\theta_v$ for every impacted edge against $H_1:\exists(u,v),\ |\theta_u-\theta_v|\ge\gamma$. This \emph{patch-neutrality} test is not the same as general thresholded preservation. Positive margin is essential: no uniformly finite worst-case query budget can distinguish exactly fair boundary instances from arbitrarily small violations under stochastic observations. The complete thresholded problem needs declared headroom or margin assumptions.

\begin{researchalgorithm}{Algorithm 1. Version-aware fairness-regression audit}\label{alg:audit}
\textbf{Input:} certified $P_0$, patch $P_1$, graph $G$, $\tau$, error $\delta$, sound impact/effect envelopes, legal replay groups.\\
\textbf{1.} Retain edges not eliminated by sound localization or $|g_{0,e}|+b_e\le\tau_e$.\\
\textbf{2.} Identify Gaussian/threshold-Bernoulli law, null scope, legal tape synchronization, cost units.\\
\textbf{3.} Choose the \emph{proved} matching schedule: $\Psi$ allocation; four-way/suite replay; $C_{\rm vert}$ raw contrast; or matched-background controls.\\
\textbf{4.} Execute the selected test; stop and decide according to its proved decision rule. Record concrete executions, error target, witness, and scope.\\
\textbf{5.} Reject unsupported composite-null or weak-replay extrapolations; return \emph{not certified} if no valid theorem applies.
\end{researchalgorithm}

\section{Gaussian paired-execution complexity}
A vertex-paired old/new difference $Z_v=P_1(v;R)-P_0(v;R)$ has mean $\theta_v$ and variance
\begin{equation}\label{eq:variance}
\nu_v^2=\mathrm{Var}(P_1(v;R))+\mathrm{Var}(P_0(v;R))-2\mathrm{Cov}(P_1(v;R),P_0(v;R)).
\end{equation}
Assume independent calls and $Z_v\sim\mathcal N(\theta_v,\nu_v^2)$ with \emph{known} $\nu_v^2>0$. Write $V_+(H)$ for vertices incident to a residual edge. One active-vertex sample supports all incident fairness edges. This reuse is captured by the convex allocation
\begin{equation}\label{eq:psi}
\Psi(H,\nu)=\min_{q_v>0,\ v\in V_+(H)}\left\{\sum_{v\in V_+(H)}q_v:\frac{\nu_u^2}{q_u}+\frac{\nu_v^2}{q_v}\le1\ (u,v)\in E(H)\right\}.
\end{equation}
\textbf{Gaussian allocation bounds.} For a nonempty residual graph, a static $\delta$-correct test takes at most $8\Psi(H,\nu)\gamma^{-2}\log(2|E(H)|/\delta)+|V_+(H)|$ Gaussian delta-oracle calls; every adaptive $\delta$-correct tester with finite null stopping expectation needs at least
\begin{equation}\label{eq:psi-lb}
\E_0[T]\ge\frac{2\kl(1-\delta,\delta)}{\gamma^2}\,\Psi(H,\nu).
\end{equation}
The proof of the upper bound scales an optimal allocation $q^\star$, thresholds empirical edge differences at $\gamma/2$, and uses Gaussian tails and a union bound. For the lower bound, change only the endpoints of each edge and optimize their mean shift under a fixed separation; adaptive transcript KL and data processing force each inequality in \Cref{eq:psi}. The results match graph/variance and margin dependence but retain a $\log|E|$ gap between the \emph{static upper} and \emph{adaptive-null lower} bounds.

The primal has a standard rotated-SOC representation, and duality gives
\begin{equation}\label{eq:psi-dual}
\Psi(H,\nu)=\max_{\lambda_e\ge0}\left[2\sum_{v\in V_+(H)}\nu_v\sqrt{\sum_{e\ni v}\lambda_e}-\sum_e\lambda_e\right].
\end{equation}
For unit variances: one edge costs $4$, a $k$-edge matching $4k$, clique $K_n$ with $n\ge2$ costs $2n$, and complete bipartite $K_{a,b}$ costs $(\sqrt a+\sqrt b)^2$. The graph is algorithmically relevant: samples at shared vertices can be reused rather than charged independently to every edge. Graph-based max--min experimental design is established \cite{SinghSinghDavidov2025}; our increment is the allocation induced by paired old/new execution, vertex reuse, and an adaptive null-side converse.

\begin{researchalgorithm}{Algorithm 2. Gaussian graph allocation}\label{alg:gaussian}
\textbf{Input:} residual $H$, known variances $\nu_v^2$, margin $\gamma$, error $\delta$.\\
If $E(H)=\varnothing$, accept with zero calls. Otherwise solve \Cref{eq:psi} for $q^\star$; for each $v\in V_+(H)$, collect $n_v=\lceil8q_v^\star\log(2|E|/\delta)/\gamma^2\rceil$ Gaussian delta observations. Reject patch-neutrality when some $|\bar Z_u-\bar Z_v|\ge\gamma/2$; otherwise accept.\\
\textbf{Guarantee:} $\delta$-correct for the stated separated Gaussian problem; explicit sample-sharing and concrete-cost accounting.
\end{researchalgorithm}

\section{Why synchronized replay changes the exponent}
For Bernoulli programs $P_j(v;U)=\mathbf 1\{U\le p_{j,v}\}$ under a common uniform tape, the signed old/new difference is supported only on an interval of length $d=|p_{1,v}-p_{0,v}|$: its variance is $d(1-d)$, whereas fresh independent sampling has constant-order variance at interior $p$. This sparse witness drives the first separation.

\textbf{One-edge replay separation.} For fixed interior $p$ and a null with identical endpoint patches, the optimal small-effect dependence under raw-pair replay is $\Theta(d^{-1}\log(1/\delta))$, with a zero-type-I witness test using
\begin{equation}\label{eq:pair-steps}
L(d,\delta)=\left\lceil\frac{\log\delta}{\log(1-d)}\right\rceil
\end{equation}
rounds. Independent new-version sampling needs a quadratic-in-$1/d$ number of observations, namely $\Theta(d^{-2}\log(1/\delta))$. The adaptive lower bound follows from binary transcript KL; the comparison assumes the specified oracle and retained baseline.

\textbf{Four-way edge replay.} Consider a fairness-neutral common shift $s>0$ and an extra differential shift $d$ on one endpoint. Sharing $U$ only across versions at one vertex fails to cancel the nuisance in the joint edge comparison; sharing $U$ across \emph{both versions and both endpoints} does cancel the common shift pathwise. For $r=1-p-s$, every adaptive $\delta$-correct four-way tester has
\begin{equation}\label{eq:four-kl}
\E_0[T]\ge\frac{\kl(1-\delta,\delta)}{r\log(r/(r-d))},
\end{equation}
whereas the sparse-witness test uses $L(d,\delta)$ rounds with zero type-I error. Ordinary vertex-paired access in this family has $\Theta(d^{-2}\log(1/\delta))$ small-$d$ complexity. This is an \emph{access-model} result, not a claim that all software can synchronize four executions.

\textbf{Tolerance phase transition.} The rare-event advantage has a sharp limit. Compare a null on the composite fairness boundary with endpoint disparity $\tau$ to an alternative of gap $\tau+d$. One synchronized round has exact null-to-alternative KL
\begin{equation}\label{eq:ktau}
K_\tau(d)=\tau\log\frac{\tau}{\tau+d}+r\log\frac r{r-d},\quad r=1-p-s-\tau.
\end{equation}
Here the first term is defined as zero at $\tau=0$. Adaptive change of measure yields a lower bound $\kl(1-\delta,\delta)/K_\tau(d)$; a fixed composite-null threshold test uses at most $\lceil8(\tau+d)\log(1/\delta)/d^2\rceil$ synchronized rounds. Under the stated interior condition, the resulting order is
\begin{equation}\label{eq:phase}
Q_{\rm thr}=\Theta\left(\frac{\tau+d}{d^2}\log\frac1\delta\right).
\end{equation}
Thus $\tau=O(d)$ retains $d^{-1}$, while fixed $\tau>0$ yields $d^{-2}$. The exact KL calculation and monotone binomial tail test are essential: this is not a universal phase diagram for arbitrary null distributions.

\begin{researchalgorithm}{Algorithm 3. Four-way replay witness}\label{alg:fourway}
\textbf{Input:} two endpoints, versions $P_0,P_1$, raw four-way common-tape access, $d$, $\delta$, and a declared homogeneous null.\\
For $t=1,\ldots,L(d,\delta)$, draw a fresh $U_t$ and evaluate both versions at both endpoints with the \emph{same} $U_t$. Reject if their old/new patch deltas differ; otherwise accept.\\
\textbf{Guarantee:} zero type-I error and pointwise type-II at most $\delta$ under the scoped one-edge family. Each round costs four concrete executions.
\end{researchalgorithm}

\textbf{Suite replay.} One shared tape across the entire active suite reveals all edge-wise patch-delta differences per round. In the homogeneous single-changed-endpoint family the tester needs $L(d,\delta)$ rounds, hence $2n_+L(d,\delta)$ concrete executions, independent of $|E|$ in \emph{round count}. The suite oracle is more informative per round but not necessarily cheaper per concrete execution.

\section{Partial replay: topology and controls}
Real systems may synchronize only subsets $S$ from a legal group family $\mathcal S$. Distinguish two incidence matrices: $A_{eS}=\mathbf1\{e\subseteq S\}$ records whether group $S$ witnesses edge $e$, whereas $B_{vS}=\mathbf1\{v\in S\}$ records whether the group exposes a changed active vertex. With group costs $c_S>0$, define
\begin{align}\label{eq:covers}
C_{\rm edge}&=\min\{c^\top a:Aa\ge\mathbf1,\ a\ge0\},\\
C_{\rm vert}&=\min\{c^\top q:Bq\ge\mathbf1,\ q\ge0\}.
\end{align}
Always $C_{\rm vert}\le C_{\rm edge}$ for the specified active set; the gap can grow. For clique $K_n$ with only pair groups and unit costs, $C_{\rm edge}=\binom n2$ and $C_{\rm vert}=n/2$, so the ratio is $n-1$. Some legal group families cover every active vertex but no group contains certain fairness-edge endpoints, making the edge-witness cover infeasible while the information cover remains finite.

\textbf{Achievability has a scope restriction.} In the homogeneous single-changed-endpoint family, a within-group vector of old/new deltas is constant under the null. Querying each group $S_i$ $n_i=\lceil L(d,\delta)q_i^\star\rceil$ times and rejecting a nonconstant vector gives zero type-I error and pointwise type-II at most $\delta$, at declared cost at most $L C_{\rm vert}+\sum_{q_i^\star>0}c_i$. Raw-transcript change of measure yields the matching active-vertex \emph{combinatorial} information floor. This extension fails for heterogeneous benign patch effects: a deterministic fair-null counterexample has new gap $0.15\le\tau=0.20$ while the homogeneous test rejects on a tape set of measure $0.15$. The more general composite null therefore remains open.

\textbf{Matched backgrounds.} A heterogeneous two-class threshold family admits another controlled repair. Some groups expose a candidate and a same-background control (first-order information); boundary-only groups do not (second-order information). With the corresponding incidence matrices $C,R$ and per-action KL values, the null expected declared cost obeys
\begin{equation}\label{eq:p19-lp}
\min_{n\ge0}\{c^\top n:(K_{\rm ctrl}C+K_{\rm bnd}R)n\ge\kl(1-\delta,\delta)\mathbf1\}.
\end{equation}
At fixed positive tolerance, $K_{\rm ctrl}=\Theta(d)$ while $K_{\rm bnd}=\Theta(d^2)$. If the fractional $C_{\rm ctrl}=\min\{c^\top q:Cq\ge\mathbf1\}$ is finite, a zero-type-I control witness uses $O(C_{\rm ctrl}d^{-1}\log(1/\delta))$ declared cost. Without a legal matched-background control, the boundary-only family remains quadratic. These statements are proved for the declared two-class/bipartite null, not arbitrary protected-group populations.

\begin{researchalgorithm}{Algorithm 4. Scope-aware partial replay}\label{alg:partial}
\textbf{Input:} legal groups, declared null family, group costs, target $d,\delta$.\\
\textbf{Homogeneous:} solve $C_{\rm vert}$; schedule $\lceil Lq_i^\star\rceil$ tapes/group and test within-group patch-delta constancy.\\
\textbf{Two-class matched-control:} solve $C_{\rm ctrl}$; schedule matched-background witness groups.\\
\textbf{Otherwise:} refuse a linear-regime guarantee; boundary-only and composite-null handling need their proper tests and validity conditions.
\end{researchalgorithm}

\textbf{Mixed-topology refinement.} For each fixed finite partial-control topology, we derive two-term small-$d$ asymptotics for its information lower program and a hybrid testing construction matching the leading \emph{orders}. A disjoint unit-cost boundary family of $m$ alternatives has a fixed nonadaptive allocation penalty $\Omega(m\log m/d^2)$, whereas sequential likelihood-ratio stopping removes the multiplicative $\log m$ from \emph{expected null cost} with $\Theta(m/d^2)$ topology scaling at fixed confidence. A deterministic truncation bound preserves the latter scale. This does \emph{not} prove exact mixed minimax constants or unrestricted optimal confidence dependence; those remain open. The proof relies on a mixture chi-square argument for fixed designs and optional-stopping-valid likelihood-ratio martingales for adaptive stopping.
Sequential multiple testing and sparse recovery already demonstrate data-dependent stopping and nonadaptive-versus-sequential dimension penalties \cite{MalekEtAl2017,MalloyNowak2014}.  Our refinement is therefore a specialized instantiation: its claimed increment is the replay-derived mixture of first-order matched-control and second-order boundary channels, not sequential removal of a multiplicity penalty in general.

\begin{table}[t]
\centering
\caption{Asymptotic testing costs in the stated Bernoulli families. These are mathematical guarantees, \emph{not measured speedups}. Different rows assume different replay oracles and null classes; group costs are not automatically concrete program executions.}
\label{tab:theory-regimes}
\normalsize
\setlength{\tabcolsep}{5pt}
\begin{tabularx}{\textwidth}{@{}p{0.32\textwidth}p{0.34\textwidth}X@{}}
\toprule
Access and hypothesis family & Established small-$d$ order & Unit and qualification\\
\midrule
Fresh independent sampling, one edge & $\Theta(d^{-2}\log(1/\delta))$ & Samples; interior known baseline\\
Paired common-uniform replay, one edge & $\Theta(d^{-1}\log(1/\delta))$ & Paired rounds; two executions each\\
Four-way replay with a benign common shift & $\Theta(d^{-1}\log(1/\delta))$ & Four-way rounds; four executions each\\
Four-way replay, fixed positive tolerance & $\Theta(d^{-2}\log(1/\delta))$ & Four-way rounds; interior boundary law\\
Matched-background controls, two-class family & $O(C_{\rm ctrl}d^{-1}\log(1/\delta))$ & Declared group-cost upper bound\\
Disjoint boundary-only channels & fixed: $\Omega(md^{-2}\log m)$; sequential: $\Theta(md^{-2})$ & Fixed confidence; latter is expected null cost\\
\bottomrule
\end{tabularx}
\end{table}

\section{Experimental verification and executable subject}
\paragraph{Evidence categories.} We distinguish (i) 821 deterministic theorem-testing specifications, (ii) seven observed population inputs in an executable software-version comparison, and (iii) stochastic simulations under the mathematical oracle models. Specification counts are not trial counts. The frozen P12/P15 campaign contains 198 configurations (54 paired-replay and 144 four-way-replay), each with 100,000 trials. Independent revalidation performed 594 exact binomial-law comparisons: none crossed the familywise $0.05$ Bonferroni threshold; the minimum exact $p$-value was $4.65\times10^{-4}$ and the maximum absolute deviation from its exact law was $2.43\times10^{-3}$. These are sufficient-statistic oracle checks, not runtime measurements, proof evidence, or evaluations of real software. Each configuration uses one theorem-prescribed sample budget, so the recorded schedule cannot independently estimate a scaling exponent. Thus \Cref{tab:theory-regimes} continues to summarize mathematical bounds rather than measured speedups.

\paragraph{Historical executable change.} A pinned historical parent/repair pair of Fairlearn's threshold optimizer \cite{BirdEtAl2020Fairlearn} provides an executable software subject. The child repairs multi-column protected-group handling; the controlled comparison holds the randomized decision law fixed. The primary metric is the exact row-wise absolute difference between the native protected-group randomized-prediction PMF and an injective canonical encoding of the same \emph{ordered tuple}. It quantifies a software-behavior representation discrepancy, \emph{not} causal discrimination or population fairness. Single-attribute populations act as negative controls. The parent differs on 244 of 4,096 Adult rows (5.96\%) and 443 of 4,096 Default Credit rows (10.82\%); the repair differs on zero rows. The full $7\times2$ parent/repair census, including five exact-zero populations, is reported in \Cref{tab:historical-census}. Source records for the two nonzero populations are \cite{AdultDatasetUCI,DefaultCreditUCI}.



These are deterministic, finite-row PMF comparisons; sampling confidence intervals do not apply. The parent defect is \emph{not} a fairness-regression-rate estimate, and its zeros are equally important: five populations do not exhibit a discrepancy for this particular exact metric. The parent/child change is a real software version transition, but it does not experimentally establish the general stochastic replay theorems. It instead validates the separate software-evolution motivation. No pretrained model or model API is a research subject.

\paragraph{Reproducibility and scope.} The analysis includes negative controls, exact fractional-cover fixtures, and a heterogeneous-null counterexample. Historical comparisons use recorded source revisions, dataset hashes, protected-category ordering, and pinned dependencies. They execute on a CPU and are distinct from the random-tape simulator. The latter's hardware throughput is not evidence of the execution cost of an actual software test. The complete P12/P15 records, hashes, seeds, manifest, and independent validation report accompany the editable artifact; no hardware speedup claim is made.

\section{Two revealing finite configurations}
\label{sec:worked}
The asymptotic theorems have different operational consequences even on small graphs. The following exact instances help separate \emph{program fairness edges}, \emph{which observations are permitted}, and \emph{which transcripts a tester can use}. They are deterministic constructions rather than empirical estimates.

\subsection{Same fairness graph, different control access}
Consider the specified bipartite two-class family with three candidates in the high-background class and one low-background reference. The three fairness edges join each high candidate to the low reference. Hold that fairness graph, the null/alternative thresholds, and all cost definitions fixed. Only the legal shared-tape groups change.

\paragraph{Boundary-only access.} Suppose the only legal groups are the three high--low pairs. Every candidate remains observable; nevertheless, each comparison puts two different background probabilities into the same tape. There is no group containing a candidate and a matched high-background control. Under a fixed positive tolerance $\tau$, the exact boundary channel obeys the small-effect expansion
\begin{equation}\label{eq:bnd-expansion}
K_{\rm bnd}(\tau,d)=\frac{d^2}{2}\left(\frac1\tau+\frac1r\right)+O(d^3),
\end{equation}
with $r$ the residual interior mass. In particular, every legal group yields only second-order information, and the necessary sample cost is quadratic in $1/d$. The software update is no less real; what disappears is the \emph{replay channel} that exposes it as a sparse witness.

\paragraph{Matched-control access.} Keep the same fairness graph but allow the three high--high pairs as replay groups. Their candidate/group incidence matrix covers the three high vertices, and symmetry yields the exact fractional cover
\begin{equation}\label{eq:triangle-cover}
\begin{aligned}
C_{\rm ctrl}&=\min_{q\ge0}(q_{12}+q_{13}+q_{23})=\tfrac32,\\
\text{s.t. }&q_{12}+q_{13}\ge1,\\
&q_{12}+q_{23}\ge1,\quad q_{13}+q_{23}\ge1.
\end{aligned}
\end{equation}
Indeed, assigning $1/2$ to each pair attains $3/2$, while summing the three constraints proves the same lower bound. Under the two-class null, a same-background control has the same pathwise old/new patch delta as an unchanged candidate. An extra shift of size $d$ at one candidate makes their deltas disagree exactly when a shared uniform tape falls in the newly added length-$d$ interval. The resulting information is first-order in $d$, and an upward-rounded fractional schedule achieves the theorem's $O(d^{-1}\log(1/\delta))$ scale. This calculation is about \emph{declared group cost}; if one wishes to compare executable cost, every two-vertex old/new group query must be converted to four program executions.

Thus the same three fairness obligations admit different asymptotic certification costs solely because the replay topology changes. The lesson is not that a favorable coupling can always be imposed, but that its availability must be part of the input specification.

\subsection{Why a witness-only planner can be suboptimal}
The gap between edge-witness coverage and raw-transcript information does not disappear by adding more pair groups. For $K_n$, allow as legal groups all size-two vertex pairs with unit cost. Every fairness edge needs its own witness group, forcing $C_{\rm edge}=\binom n2$. But every vertex participates in $n-1$ groups; assigning weight $1/(n-1)$ to each group meets the active-vertex constraints, yielding $C_{\rm vert}=n/2$. Their ratio is exactly $n-1$, unbounded as $n$ grows.

A stronger obstruction uses the three-vertex path with legal groups $\{0,2\}$ and $\{1,2\}$. The edge $\{0,1\}$ is in no legal group, hence $C_{\rm edge}=\infty$; nevertheless, all relevant active vertices appear in at least one group, and $C_{\rm vert}$ is finite. Within the specified homogeneous single-changed-endpoint family, full raw old/new output vectors make it possible to exploit that latter information. Under arbitrary heterogeneous fair nulls, the same raw-contrast predicate can fire spuriously, so LP feasibility alone does not establish validity for composite nulls.

\subsection{What the experiments can and cannot establish}
The theoretical instances fix threshold distributions and permit exact KL and fractional-LP evaluation. Their simulations can falsify an implementation by comparing finite-sample outcomes to these laws, but agreement does not prove theorems or demonstrate real-world causal fairness. In contrast, the historical protected-tuple defect tests whether a genuine old/new software change disagrees with a canonical group representation. The two positive parent-revision comparisons and five negative controls provide source-level evidence without pretending that the finite census estimates an underlying population effect. Together, these two evidence layers connect mathematical access models to real evolving code while preserving the distinction between the abstractions and the observed software behavior.

\section{Relation to prior work and limits}
Software fairness testing and directed test generation motivate the specification but do not provide our replay hierarchy \cite{GalhotraBrunMeliou2017,UdeshiAroraChattopadhyay2018,ZhangEtAl2020,ChenEtAl2024}. Yan and Zhang actively audit demographic parity for a fixed model \cite{YanZhang2022}; Ajarra and Basu study PAC auditing after strategic model updates via their SP dimension for statistical parity \cite{AjarraBasu2026}. Neither work models executable old/new transcript coupling or charges concrete executions. The method in \cite{BuselliEtAl2024} mitigates group imbalance in negative prediction flips after model updates, whereas our edgewise property and query-complexity objective are different. Incremental and regression testing explain reuse across versions \cite{RothermelHarrold1996,YooHarman2012,LiSunLeungZhang2013,YinEtAl2025}. Fixed-model individual-fairness verification and relational $k$-safety specifications \cite{GeorgeJohnVijaykeerthySaha2020,ChristakisEtAl2023} do not by themselves establish version-aware stochastic replay bounds. Efficient action selection for active sequential testing is likewise established \cite{GanJiaLiTayur2026}; our claimed distinctions therefore concern the exact legal old/new transcript, fairness-null family, and concrete execution accounting, not sequential sensing as such. Common random numbers and paired designs are standard \cite{GlassermanYao1992,Connor1987,McNemar1947,NelsonMatejcik1995}; adaptive allocation, sequential multiple testing, and controlled sensing already have deep theory \cite{Chernoff1959,CohenZhao2015,NaghshvarJavidi2013,NitinawaratAtiaVeeravalli2013,DeshmukhVeeravalliBhashyam2021,MalekEtAl2017,MalloyNowak2014}; graph-structured design and property testing supply related frameworks \cite{SinghSinghDavidov2025,GoldreichGoldwasserRon1998,GoldreichRon1997,BlaisSeth2024,BunDesaiPinto2026,PintoPalitRaskhodnikova2026}.

Our distinct claim is a \emph{version-aware, fairness-graph-constrained} hierarchy of legal replay transcripts and costs, where explicit null classes permit both exponent shifts and converse results. The study does not characterize general composite nuisance classes (a general composite-nuisance characterization remains conjectural), unknown-variance anytime-valid procedures, arbitrary unsound impact reports, uncontrolled software nondeterminism, or causal discrimination in the observational populations. These distinctions are essential to separating the present statements from classical controlled sensing, coupling, and graph-structured experimental design.

\section{Validity conditions and limitations}
\label{sec:gate}
\paragraph{Which claims have matching hypotheses?}
The Gaussian upper and lower bounds hold with independent calls, known positive vertex variances, a positive separation margin, and the specified delta oracle. The Bernoulli $d^{-1}$ advantage assumes a common threshold tape and a sparse changed interval. Suite-replay results assume the entire active suite can share the same tape; the execution cost grows with the number of executed vertices. Partial-replay $C_{\rm vert}$ achievability uses a homogeneous single-changed-endpoint null, and the heterogeneous control result uses declared control classes. None of these assumptions follows merely from the existence of a protected attribute or from GPU simulation.

\paragraph{Avoiding incorrect lower-bound transference.}
A transcript-KL lower bound depends on what the observation reveals. Replacing a full raw transcript with only an edgewise witness can make an apparent lower bound too large: the partial-replay information geometry is governed by active vertices, not edges, in its homogeneous family. Similarly, for a composite null there is one unknown but globally fixed nuisance parameter; minimizing allocations separately at each null parameter and imposing a common allocation against all nulls lead to different mathematical quantities. We do not promote either interchange into a general adaptive minimax equivalence. A matching minimax theorem for this broader composite-nuisance setting remains open.

\paragraph{Unit comparability and adaptivity.}
A vertex-paired physical execution uses two program executions whether it exposes a raw pair or only a Gaussian delta; an edge-synchronized call uses four; a complete suite round costs $2n_+$ executions. Group costs may represent $2|S|$ runs or another explicitly declared charge. A theorem comparing $d$-exponents for different oracle classes does not automatically imply wall-clock or energy savings. The familywise $\log m$ penalty applies to deterministic fixed boundary allocations in the stated independent-channel family; sequential stopping improves expected null cost. The resulting bound is neither a uniform improvement for all alternative stopping times nor proof of exact optimal constants in mixed topologies.

\paragraph{Limits of interpretation.}
The theory provides algorithmic and information-theoretic certificates for explicit program laws. It does not establish fairness of every updated deployed service, soundness of arbitrary static analysis, replay of nonseedable external nondeterminism, or freedom from protected-feature leakage in an unrelated model pipeline. The central value is sharper \emph{verification cost accounting}, not a universal method for measuring societal fairness. The role of replay topology and certificate reuse must therefore be interpreted alongside prior results in controlled sensing, adaptive testing, coupling, regression selection, and property testing.

\noindent\textit{The next section states the principal results formally and supplies proof arguments; additional supporting proofs and experimental details appear in the appendix.}\par
\medskip
\section{Formal results and proof architecture}
\label{sec:formal}
The preceding sections provide an accessible map of the results. We now record formal statements and principal proofs, with explicit quantifiers and risk guarantees. These formulas should be evaluated within their declared observation families; an asymptotic law derived for one family is not a generic fairness-auditing guarantee.

\paragraph{Observation primitives.}
A Gaussian delta-oracle call reveals only
$Z_v=P_1(v;R)-P_0(v;R)$.  A raw vertex-pair replay reveals the ordered pair
$(P_0(v;R),P_1(v;R))$.  Raw edge, suite, and legal-group replay reveal all
ordered old/new outputs generated with their declared common tape.  Different
rounds use independent tapes.  These primitives can have the same physical
execution count but different sigma-algebras and are not interchangeable in a
converse.  Write $V_+(H)=\{v:\deg_H(v)>0\}$.


\subsection{The Gaussian allocation converse and primal design}


\begin{theorem}[Static upper bound]
\label{thm:upper}
Let $m=|E(H)|\ge1$.  For $\delta\in(0,1)$, a non-adaptive Gaussian
delta-oracle tester exists using at
most
\begin{equation}\label{eq:formal-results-001}
\frac{8\log(2m/\delta)}{\gamma^2}\Psi(H,\nu)+|V_+(H)|
\end{equation}
Gaussian delta-oracle calls and having error probability at most $\delta$.
\end{theorem}

\begin{proof}
Take an optimal solution $q^\star$ and set
\begin{equation}\label{eq:formal-results-002}
n_v=\left\lceil
\frac{8\log(2m/\delta)}{\gamma^2}q_v^\star
\right\rceil.
\end{equation}
For each edge $(u,v)$, the empirical difference of the two vertex means has
variance at most
\begin{equation}\label{eq:formal-results-003}
\frac{\gamma^2}{8\log(2m/\delta)}.
\end{equation}
Reject if an empirical edge
difference has magnitude at least $\gamma/2$.  Under the null, a Gaussian tail
bound followed by a union bound controls all edges.  Under the alternative, an
edge with true difference at least $\gamma$ can be missed only if its estimation
error is at least $\gamma/2$.  The ceilings contribute at most $|V_+(H)|$
Gaussian delta-oracle calls.
\end{proof}


\begin{theorem}[Adaptive lower bound]
\label{thm:lower}
For $\delta\in(0,1/2)$, every adaptive $\delta$-correct tester with finite
expected stopping time (measured in Gaussian delta-oracle calls) has, on the null instance $\theta\equiv0$,
\begin{equation}\label{eq:formal-results-004}
\E_0[T]
\ge
\frac{2\kl(1-\delta,\delta)}{\gamma^2}\Psi(H,\nu).
\end{equation}
\end{theorem}

\begin{proof}
Let $n_v=\E_0[N_v]$.  Fix an edge $(u,v)$ and consider an alternative changing
only its endpoints to $\theta_u=a$ and $\theta_v=-b$, where $a,b\ge0$ and
$a+b=\gamma$.  The adaptive transcript KL from the null to this alternative is
\begin{equation}\label{eq:formal-results-005}
\frac{n_u a^2}{2\nu_u^2}+\frac{n_v b^2}{2\nu_v^2}.
\end{equation}
The standard adaptive change-of-measure argument for sequential sampling
\cite{Chernoff1959,CohenZhao2015}, followed by data processing through the final
accept/reject decision, lower-bounds this by $\kl(1-\delta,\delta)$.  Minimizing over $a+b=\gamma$ yields
\begin{equation}\label{eq:formal-results-006}
\frac{\nu_u^2}{n_u}+\frac{\nu_v^2}{n_v}
\le
\frac{\gamma^2}{2\kl(1-\delta,\delta)}.
\end{equation}
This holds for every edge.  Scaling $n_v$ by
$\gamma^2/(2\kl(1-\delta,\delta))$ therefore gives a feasible solution to
\Cref{eq:psi}, proving the claim.
\end{proof}


\subsection{Replay versus fresh observation}


For the one-edge family, the old thresholds are $p$ at both endpoints. The null applies no differential patch and the alternative adds $d$ to the first endpoint; in four-way replay a common shift $s$ is present in both new endpoint thresholds. These access-specific constructions supply the reference null and alternative distributions used below.


\begin{theorem}[Replay changes small-patch complexity]
\label{thm:bernoulli-separation}
Fix $p\in(0,1)$, $d\in(0,1-p)$, and $\delta\in(0,1/2)$.  For the one-edge family above:

\begin{enumerate}[label=(\roman*)]
\item Under common-uniform raw-pair replay, there is a test using
\begin{equation}\label{eq:formal-results-007}
n_{\rm pair}
=
\left\lceil
\frac{\log \delta}{\log(1-d)}
\right\rceil
\end{equation}
paired rounds with zero type-I error and type-II error at most $\delta$.
Moreover, every adaptive $\delta$-correct paired-replay tester with finite
expected stopping time satisfies, on the null,
\begin{equation}\label{eq:formal-results-008}
\E_0[T]
\ge
\frac{\kl(1-\delta,\delta)}
{(1-p)\log\!\left(\frac{1-p}{1-p-d}\right)}.
\end{equation}

\item Under independent new-version observation, a test using
\begin{equation}\label{eq:formal-results-009}
n_{\rm ind}
=
\left\lceil
\frac{2\log(1/\delta)}{d^2}
\right\rceil
\end{equation}
single-execution new-version calls at $u$ is $\delta$-correct when the retained old baseline
$p$ is known.  Conversely, every adaptive $\delta$-correct independent-replay
tester with finite expected stopping time satisfies
\begin{equation}\label{eq:formal-results-010}
\E_0[T]
\ge
\frac{\kl(1-\delta,\delta)}
{\kl(p,p+d)}.
\end{equation}
\end{enumerate}

Consequently, for every fixed interior $p$,
\begin{equation}\label{eq:formal-results-011}
Q_{\rm pair}(d,\delta)
=
\Theta\!\left(
\frac{\log(1/\delta)}{d}
\right),
\qquad
Q_{\rm ind}(d,\delta)
=
\Theta\!\left(
\frac{\log(1/\delta)}{d^2}
\right)
\end{equation}
as $d\to0$ with $p$ and $\delta$ fixed, up to constants depending on them.
\end{theorem}

\begin{proof}
For paired replay, run $P_0$ and $P_1$ on $u$ using the same
$U\sim\mathrm{Unif}[0,1]$ and reject upon observing a discordant pair
$(Y_0,Y_1)=(0,1)$.  Under $\mathsf H_0$ discordance is impossible.  Under
$\mathsf H_1$ it occurs exactly when $U\in(p,p+d]$, hence independently with
probability $d$ each round.  The miss probability after $n$ rounds is therefore
$(1-d)^n$, proving the upper bound.

For the paired lower bound, one informative query at $u$ has null law
\begin{equation}\label{eq:formal-results-012}
Q_0(1,1)=p,
\qquad
Q_0(0,0)=1-p,
\end{equation}
and alternative law
\begin{equation}\label{eq:formal-results-013}
Q_1(1,1)=p,
\qquad
Q_1(0,0)=1-p-d,
\qquad
Q_1(0,1)=d.
\end{equation}
Thus
\begin{equation}\label{eq:formal-results-014}
\KL(Q_0\|Q_1)
=
(1-p)\log\!\left(\frac{1-p}{1-p-d}\right).
\end{equation}
Queries to $v$ have identical laws in both hypotheses.  The adaptive transcript
chain rule and data processing through the terminal accept/reject decision give
\begin{equation}\label{eq:formal-results-015}
\E_0[N_u]\KL(Q_0\|Q_1)
\ge
\kl(1-\delta,\delta),
\end{equation}
which yields the stated lower bound.

For independent replay, giving the tester the exact old baseline $p$ can only make
the problem easier.  Sampling the new version at $u$ and rejecting when its
empirical mean is at least $p+d/2$ has each one-sided error bounded by
$\exp(-nd^2/2)$ by Hoeffding's inequality, proving the upper bound.
For the lower bound, the old-version observation has the same distribution under
both hypotheses and contributes no information.  Each informative new-version
sample contributes
\begin{equation}\label{eq:formal-results-016}
\KL(\mathrm{Bern}(p)\|\mathrm{Bern}(p+d))
=
\kl(p,p+d).
\end{equation}
The same transcript change-of-measure argument yields the claimed expected-sample
lower bound.

Finally,
\begin{equation}\label{eq:formal-results-017}
(1-p)\log\!\left(\frac{1-p}{1-p-d}\right)
=
d+O(d^2),
\end{equation}
whereas
\begin{equation}\label{eq:formal-results-018}
\kl(p,p+d)
=
\frac{d^2}{2p(1-p)}+O(d^3).
\end{equation}
The asymptotic separation follows.
\end{proof}


\begin{theorem}[Counterfactual replay changes the small-gap exponent]
\label{thm:counterfactual-replay}
Fix \(p>0\), \(s>0\), \(d>0\), \(p+s+d<1\), and
\(\delta\in(0,1/2)\).  The old thresholds are \((p,p)\); the null new
thresholds are \((p+s,p+s)\), and the alternative new thresholds are
\((p+s+d,p+s)\).  Vertex-paired and four-way rounds reveal their ordered raw
outputs, and tapes are independent across rounds.

\begin{enumerate}[label=(\roman*)]
\item Under four-way counterfactual replay, there is a test using
\begin{equation}\label{eq:formal-results-019}
n_{\rm edge}
=
\left\lceil
\frac{\log\delta}{\log(1-d)}
\right\rceil
\end{equation}
replay rounds with zero type-I error and type-II error at most \(\delta\).
Every adaptive \(\delta\)-correct tester with finite expected stopping time
satisfies, on the null,
\begin{equation}\label{eq:formal-results-020}
\E_0[T]
\ge
\frac{\kl(1-\delta,\delta)}
{r\log\!\left(\frac{r}{r-d}\right)},
\qquad
r=1-p-s.
\end{equation}

\item Under vertex-paired replay, even granting the tester \(p\), \(s\), and the
identity of the endpoint that may receive the differential effect, a
\(\delta\)-correct test exists using
\begin{equation}\label{eq:formal-results-021}
n_{\rm vertex}
=
\left\lceil
\frac{2\log(1/\delta)}{d^2}
\right\rceil
\end{equation}
raw vertex-pair rounds.  Conversely, every adaptive \(\delta\)-correct
vertex-paired tester with finite expected stopping time satisfies
\begin{equation}\label{eq:formal-results-022}
\E_0[T]
\ge
\frac{\kl(1-\delta,\delta)}
{
s\log\!\left(\frac{s}{s+d}\right)
+
r\log\!\left(\frac{r}{r-d}\right)
}.
\end{equation}
\end{enumerate}

Consequently, for fixed \(p,s,r>0\),
\begin{equation}\label{eq:formal-results-023}
Q_{\rm edge}(d,\delta)
=
\Theta\!\left(
\frac{\log(1/\delta)}{d}
\right),
\qquad
Q_{\rm vertex}(d,\delta)
=
\Theta\!\left(
\frac{\log(1/\delta)}{d^2}
\right)
\end{equation}
as \(d\to0\) with \(\delta\) fixed, up to constant factors.  Counting concrete program executions
multiplies the two costs by four and two, respectively, and therefore does not
change the exponent separation.
\end{theorem}

\begin{proof}
Under four-way replay, define the pathwise differential
\begin{equation}\label{eq:formal-results-024}
W
=
\bigl(P_1(u;U)-P_0(u;U)\bigr)
-
\bigl(P_1(v;U)-P_0(v;U)\bigr).
\end{equation}
Under \(\mathsf H_0\), the two endpoint deltas are identical for every \(U\), so
\(W=0\) almost surely.  Under \(\mathsf H_1\),
\begin{equation}\label{eq:formal-results-025}
W=1
\quad\Longleftrightarrow\quad
U\in(p+s,p+s+d],
\end{equation}
which occurs with probability \(d\).  Rejecting upon the first \(W=1\) therefore
has zero type-I error and miss probability \((1-d)^n\), proving the upper bound.

For the lower bound, the four-output observation has the same mass \(p\) on the
all-one state and the same mass \(s\) on the state in which both new-version
outputs are one while both old-version outputs are zero.  Under the null, the
remaining all-zero state has mass \(r=1-p-s\).  Under the alternative, mass \(d\)
is split from that state into a new state witnessing the counterfactual
differential.  Hence the null-to-alternative KL of one four-way round is
\begin{equation}\label{eq:formal-results-026}
K_{\rm edge}
=
r\log\!\left(\frac{r}{r-d}\right).
\end{equation}
Adaptive transcript change of measure followed by data processing through the
terminal decision gives
\begin{equation}\label{eq:formal-results-027}
\E_0[T]K_{\rm edge}
\ge
\kl(1-\delta,\delta).
\end{equation}

For vertex-paired replay, querying the potentially changed endpoint \(u\) produces
the pair distribution
\begin{equation}\label{eq:formal-results-028}
(p,\ s,\ r)
\end{equation}
on states \((1,1),(0,1),(0,0)\) under the null, and
\begin{equation}\label{eq:formal-results-029}
(p,\ s+d,\ r-d)
\end{equation}
under the alternative.  Its KL is therefore
\begin{equation}\label{eq:formal-results-030}
K_{\rm vertex}
=
s\log\!\left(\frac{s}{s+d}\right)
+
r\log\!\left(\frac{r}{r-d}\right).
\end{equation}
Queries to \(v\) carry no information, so the same transcript argument proves the
lower bound.  For the upper bound, ignore the extra pair information and retain
only the discordance indicator.  It is Bernoulli\((s)\) under the null and
Bernoulli\((s+d)\) under the alternative.  Thresholding its empirical mean at
\(s+d/2\) and applying Hoeffding's inequality bounds each one-sided error by
\(\exp(-nd^2/2)\).

Finally,
\begin{equation}\label{eq:formal-results-031}
K_{\rm edge}
=
d+O(d^2),
\end{equation}
while the first-order terms in \(K_{\rm vertex}\) cancel and
\begin{equation}\label{eq:formal-results-032}
K_{\rm vertex}
=
\frac{d^2}{2}
\left(
\frac{1}{s}+\frac{1}{r}
\right)
+
O(d^3).
\end{equation}
The claimed asymptotic rates follow.
\end{proof}


\subsection{Threshold nulls and suite synchronization}

For P18, fix $p>0$, $s\ge0$, $\tau\ge0$, and $d>0$ with
$p+s+\tau+d<1$.  The old thresholds are $(p,p)$; the boundary-null new
thresholds are $(p+s+\tau,p+s)$ and the alternative thresholds are
$(p+s+\tau+d,p+s)$.  Each round reveals the ordered raw four-way transcript.


\begin{theorem}[Thresholded replay phase transition]
\label{thm:threshold-replay-phase}
Let $\delta\in(0,1/2)$ and set $r=1-p-s-\tau$.  One four-way synchronized
replay round has null-to-alternative KL
\begin{equation}\label{eq:formal-results-033}
K_\tau(d)
=
\tau\log\!\left(\frac{\tau}{\tau+d}\right)
+
r\log\!\left(\frac{r}{r-d}\right),
\end{equation}
where the first term is zero when $\tau=0$.  Every adaptive
$\delta$-correct tester with finite null expected stopping time satisfies
\begin{equation}\label{eq:formal-results-034}
\E_0[T]\ge
\frac{\kl(1-\delta,\delta)}{K_\tau(d)}.
\end{equation}
Conversely, the full composite problem
\begin{equation}\label{eq:formal-results-035}
\mathsf H_0:\ |p_{1,u}-p_{1,v}|\le\tau,
\qquad
\mathsf H_1:\ |p_{1,u}-p_{1,v}|\ge\tau+d
\end{equation}
admits a fixed-design synchronized-replay test using at most
\begin{equation}\label{eq:formal-results-036}
\left\lceil
\frac{8(\tau+d)\log(1/\delta)}{d^2}
\right\rceil
\end{equation}
rounds.  If $r-d$ is bounded below by a positive constant, then as $d\to0$
with $\delta$ fixed these bounds give
\begin{equation}\label{eq:formal-results-037}
Q_{\rm thr}(\tau,d,\delta)
=
\Theta\!\left(
\frac{\tau+d}{d^2}\log\frac1\delta
\right)
\end{equation}
up to constants.  Hence $\tau=0$ and, more generally, $\tau=O(d)$ retain
$d^{-1}$ scaling, whereas every fixed $\tau>0$ restores $d^{-2}$ scaling as
$d\to0$.
\end{theorem}

\begin{proof}
Under the boundary pair, only two four-way transcript atoms change probability:
masses $\tau$ and $r$ under the null become $\tau+d$ and $r-d$ under the
alternative, giving the displayed KL.  Adaptive transcript change of measure
and data processing yield the lower bound.  Equivalently, conditioning on the
event $U>p+s$ reduces the informative transcript exactly to a Bernoulli shift
from $\tau/(1-p-s)$ to $(\tau+d)/(1-p-s)$.

For the upper bound, in each independent synchronized round retain only
\begin{equation}\label{eq:formal-results-038}
W=\mathbf 1\{P_1(u;U)\ne P_1(v;U)\}.
\end{equation}
Under common-uniform threshold replay, $W$ is Bernoulli with mean equal to the
new-version endpoint gap.  Reject when its empirical mean is at least
$\tau+d/2$.  Binomial stochastic ordering reduces the two composite errors to
the boundary means $\tau$ and $\tau+d$; Bernstein/Chernoff bounds make both at
most $\delta$ at the stated sample size.  Finally, for fixed positive $\tau$
and $r$,
\begin{equation}\label{eq:formal-results-039}
K_\tau(d)
=
\frac{d^2}{2}\left(\frac1\tau+\frac1r\right)+O(d^3),
\end{equation}
whereas $K_0(d)=d+O(d^2)$.  More uniformly, if $r-d\ge\rho>0$,
$K_\tau(d)=O_\rho(d^2/(\tau+d))$, matching the upper order.
\end{proof}


\begin{theorem}[Synchronized suite replay]
\label{thm:suite-replay}
Let \(H=(V,E)\) contain at least one fairness edge and fix
\(p,s,d>0\) with \(p+s+d<1\) and \(\delta\in(0,1/2)\).  Under the null,
every active vertex has old threshold \(p\) and new threshold \(p+s\); under
an alternative, exactly one endpoint of an active edge has new threshold
\(p+s+d\). A raw suite round
reveals both version outputs at every active vertex under one common tape.
There is a test that distinguishes the null from the union of all
single-changed-endpoint alternatives using
\begin{equation}\label{eq:formal-results-040}
T_{\rm suite}
=
\left\lceil
\frac{\log\delta}{\log(1-d)}
\right\rceil
\end{equation}
synchronized suite-replay rounds, with zero type-I error and type-II error at
most \(\delta\) for every alternative.  Hence its concrete execution cost is
\begin{equation}\label{eq:formal-results-041}
2n_+
\left\lceil
\frac{\log\delta}{\log(1-d)}
\right\rceil.
\end{equation}
Moreover, every adaptive \(\delta\)-correct tester whose primitive observation is
a synchronized suite-replay round and whose stopping time has finite expectation
satisfies, on the null,
\begin{equation}\label{eq:formal-results-042}
\E_0[T]
\ge
\frac{\kl(1-\delta,\delta)}
{r\log\!\left(\frac{r}{r-d}\right)},
\qquad
r=1-p-s.
\end{equation}
Consequently, for fixed interior \(p,s,r\),
\begin{equation}\label{eq:formal-results-043}
Q_{\rm suite}^{\rm rounds}
=
\Theta\!\left(\frac{\log(1/\delta)}{d}\right),
\qquad
Q_{\rm suite}^{\rm exec}
=
\Theta\!\left(
\frac{n_+\log(1/\delta)}{d}
\right)
\end{equation}
as \(d\to0\) with \(\delta\), the graph, and hence \(n_+\) fixed.  The round bound has no explicit \(|E|\) factor; the concrete
execution cost remains $2n_+$ per round and can scale with the graph.
\end{theorem}

\begin{proof}
In one suite round, form the old/new delta at every active vertex and compare
those deltas across every fairness edge.  Under the null, all vertex deltas are
identical pathwise because every vertex uses the same \(U\), old threshold
\(p\), and new threshold \(p+s\).  Therefore the tester never rejects.

Under an alternative, let \(u\) be the endpoint receiving the additional shift
\(d\).  Whenever
\begin{equation}\label{eq:formal-results-044}
U\in(p+s,p+s+d],
\end{equation}
the delta at \(u\) equals one while every unchanged neighboring endpoint has
delta zero.  At least one incident fairness edge therefore witnesses a
differential.  This event has probability \(d\) in every suite round, so the
probability of missing it for \(T\) rounds is \((1-d)^T\).  The stated upper
bound follows.

For the lower bound, fix any one legal alternative.  Under the null, the suite
observation has an all-common state of mass \(r=1-p-s\).  Under the fixed
alternative, mass \(d\) is split from this state into a new state identifying the
differential endpoint; all other masses are unchanged.  Hence the
null-to-alternative KL of one suite round is
\begin{equation}\label{eq:formal-results-045}
K_{\rm suite}
=
r\log\!\left(\frac{r}{r-d}\right).
\end{equation}
The standard adaptive transcript change-of-measure argument and data processing
through the terminal decision give
\begin{equation}\label{eq:formal-results-046}
\E_0[T]K_{\rm suite}
\ge
\kl(1-\delta,\delta).
\end{equation}
Finally \(K_{\rm suite}=d+O(d^2)\), proving the asymptotic statement.  Each suite
round performs two version executions at each of the \(n_+\) active inputs.
\end{proof}


\subsection{Partial replay structural and achievable geometry}


\begin{proposition}[Separation of the P17 LP geometries]
\label{prop:p17-geometry-separation}
The two fractional-cover geometries need not be comparable within any finite
universal factor.  More precisely: (i) if every fairness edge has endpoints
with identical legal-group incidence signatures, then
$C_{\rm vert}=C_{\rm edge}$; (ii) for $H=K_n$ with the legal groups equal to
its edges and unit group costs,
$C_{\rm vert}=n/2$ and $C_{\rm edge}=\binom n2$, so their ratio is $n-1$;
and (iii) there is a legal partial-replay instance with
$C_{\rm vert}<\infty$ but $C_{\rm edge}=\infty$.
\end{proposition}

\begin{proof}
For (i), identical endpoint incidence makes the edge-witness row
$\mathbf 1\{u,v\in S_i\}_i$ identical to either endpoint's active-vertex
row; equality propagates over each connected fairness component, and duplicate
constraints may be removed.  For (ii), pair groups cover only their own edge, so
every edge variable in the witness LP is at least one.  Summing the $n$
active-vertex constraints counts every pair variable twice and gives
$C_{\rm vert}\ge n/2$, while assigning weight $1/(n-1)$ to every pair
attains equality.  For (iii), take the path $0$--$1$--$2$ with legal groups
$\{0,2\}$ and $\{1,2\}$ at unit cost.  The active-vertex LP has value $2$,
but edge $\{0,1\}$ belongs to no legal group, so the witness LP is infeasible.
\end{proof}


\begin{proposition}[Active-vertex raw-contrast schedule]
\label{prop:p17-active-contrast}
Assume every vertex appearing in a legal replay group has old threshold $p$
and, under the null, new threshold $p+s$.  Under each alternative, exactly one
fairness-active endpoint $x$ instead has new threshold $p+s+d$, with
$p,s,d>0$ and $p+s+d<1$.  Assume also that every legal group in this
proposition contains at least two distinct vertices.  Suppose
$C_{\rm vert}<\infty$.  Let $q^\star$
solve the active-vertex LP, set
\begin{equation}\label{eq:formal-results-047}
L=\left\lceil\frac{\log\delta}{\log(1-d)}\right\rceil,
\qquad
n_i=\lceil Lq_i^\star\rceil,
\end{equation}
and query $S_i$ exactly $n_i$ times with independent tapes across queries.
Within a query, define $D_v=P_1(v;U)-P_0(v;U)$ and reject if the vector
$(D_v:v\in S_i)$ is nonconstant.  The test has zero type-I error and
pointwise type-II error at most $\delta$, with declared cost at most
\begin{equation}\label{eq:formal-results-048}
L C_{\rm vert}+\sum_{i:q_i^\star>0}c_i.
\end{equation}
Consequently the corrected information lower bound and this executable upper
construction use the same active-vertex combinatorial geometry in the specified
homogeneous single-changed-endpoint family.
\end{proposition}

\begin{proof}
Under the null, common-uniform replay gives
$D_v=\mathbf 1\{U\in(p,p+s]\}$ for every member of every legal group, so
the within-group delta vector is constant pathwise and rejection is impossible.
Under an alternative changing $x$, any legal group containing $x$ contains at
least one distinct member $y$ because legal groups have size at least two.
The alternative changes only $x$, hence
\begin{equation}\label{eq:formal-results-049}
D_x\ne D_y
\quad\Longleftrightarrow\quad
U\in(p+s,p+s+d],
\end{equation}
which has probability $d$.  For every active $x$, feasibility of $q^\star$
and upward rounding give
\begin{equation}\label{eq:formal-results-050}
\sum_{i:x\in S_i}n_i
\ge L\sum_{i:x\in S_i}q_i^\star
\ge L.
\end{equation}
Independent tapes across group queries therefore make the miss probability for
alternative $x$ at most $(1-d)^L\le\delta$.  Finally,
$\lceil Lq_i^\star\rceil\le Lq_i^\star+1$ on the positive support,
which gives the stated cost bound.
\end{proof}


\subsection{Heterogeneous matched-control information}


Here $C$ records which legal groups provide same-background control observations and $R$ records boundary-only observations, with positive declared group costs $c$. Every legal group has at least two vertices and reveals its ordered raw old/new transcript under one common tape; rounds are independent. Fix $p,s,\tau,d>0$ with $p+s+\tau+d<1$. Under the simple null all old thresholds equal $p$, low new thresholds equal $p+s$, and high new thresholds equal $p+s+\tau$; alternative $x$ raises only high candidate $x$ by $d$. The following statements do not cover unrestricted heterogeneous populations or simultaneous changes.


\begin{theorem}[Matched controls under a heterogeneous fair null]
\label{thm:control-replay}
Fix \(\delta\in(0,1/2)\) and the declared two-class family above.

\begin{enumerate}[label=(\roman*)]
\item Every adaptive \(\delta\)-correct partial-replay tester with finite null
expected declared cost satisfies
\begin{equation}\label{eq:formal-results-051}
\E_0[\mathrm{cost}]
\ge
\min_{n\ge0}
\left\{
c^\top n:
\bigl(K_{\rm ctrl}C+K_{\rm bnd}R\bigr)n
\ge
\kl(1-\delta,\delta)\mathbf 1
\right\}.
\end{equation}

\item Define the fractional control-cover quantity
\begin{equation}\label{eq:formal-results-052}
C_{\rm ctrl}
=
\min_{q\ge0}
\left\{
c^\top q:
Cq\ge\mathbf 1
\right\}.
\end{equation}
If \(C_{\rm ctrl}<\infty\), a fixed-design tester has zero type-I error and
pointwise type-II error at most \(\delta\) using declared cost at most
\begin{equation}\label{eq:formal-results-053}
L\,C_{\rm ctrl}
+
\sum_{i:q_i^\star>0}c_i,
\qquad
L=
\left\lceil
\frac{\log\delta}{\log(1-d)}
\right\rceil.
\end{equation}

\item For fixed \(p,s,\tau>0\) with \(p+s+\tau<1\),
\begin{equation}\label{eq:formal-results-054}
K_{\rm ctrl}(d)=d+O(d^2),
\qquad
K_{\rm bnd}(\tau,d)
=
\frac{d^2}{2}
\left(
\frac1\tau+\frac1r
\right)
+O(d^3).
\end{equation}
Consequently, if \(C_{\rm ctrl}<\infty\), the lower program and upper
construction have matching leading control-cover geometry and
\(\Theta(d^{-1}\log(1/\delta))\) fixed-confidence order, up to rounding and
standard constant factors.  If some observable high candidate has no
control-compatible legal group, then every tester requires
\(\Omega(d^{-2}\log(1/\delta))\) expected declared cost at fixed
\(\tau>0\).  If a candidate is completely unobservable, no finite-cost
\(\delta\)-correct tester exists.

\item The mixed-topology information floor has an exact two-term singular
expansion.  Let
\begin{equation}\label{eq:formal-results-055}
\begin{aligned}
U&=\{x\in A:C_{xi}=0\ \forall i\},\\
b_2&=\tfrac12(\tau^{-1}+r^{-1}),\\
b_3&=\tfrac13(r^{-2}-\tau^{-2}).
\end{aligned}
\end{equation}
and, for each observable \(x\in U\), let
\(\bar c_x=\min\{c_i:x\in S_i\}\).  Write
\begin{equation}\label{eq:formal-results-056}
C_{\rm ctrl}^{A\setminus U}
=
\min_{q\ge0}
\left\{
c^\top q:C_{A\setminus U}q\ge\mathbf 1
\right\},
\end{equation}
with value zero if \(A\setminus U\) is empty, and let
\(a=\kl(1-\delta,\delta)\).  If every candidate is observable, then as
\(d\downarrow0\) the value \(Q_{\rm LB}(d)\) of the lower program in (i)
satisfies
\begin{equation}\label{eq:formal-results-057}
Q_{\rm LB}(d)
=
\frac{a}{b_2}
\left(\sum_{x\in U}\bar c_x\right)d^{-2}
+
a\left(
C_{\rm ctrl}^{A\setminus U}
-\frac{b_3}{b_2^2}\sum_{x\in U}\bar c_x
\right)d^{-1}
+O(1).
\end{equation}
This is an exact asymptotic for the information lower program, not a matching
upper-bound expansion when \(U\ne\varnothing\).

\item Suppose \(U\ne\varnothing\) and every candidate is observable, and let
\(m=|U|\).  There is a fixed-design hybrid tester with type-I error and every
pointwise type-II error at most \(\delta\), and declared cost at most
\begin{equation}\label{eq:formal-results-058}
N_U\sum_{x\in U}\bar c_x
+L C_{\rm ctrl}^{A\setminus U}
+\sum_{i:q_i^\star>0}c_i,
\end{equation}
where
\begin{equation}\label{eq:formal-results-059}
N_U=
\left\lceil
\frac{8(\tau+d)}{d^2}\log\frac{m}{\delta}
\right\rceil,
\qquad
L=\left\lceil\frac{\log\delta}{\log(1-d)}\right\rceil,
\end{equation}
and \(q^\star\) is an optimizer defining
\(C_{\rm ctrl}^{A\setminus U}\).  Consequently, for every fixed finite
topology and fixed \(\tau>0\), the mixed lower and upper bounds match the
\(d^{-2}\) leading order and the boundary-minimum/control-cover cost geometry
up to constants.  This does not match the exact two-term coefficients or close
the dependence on \(m\) within the broader adaptive expected-cost class; the
next part separates those regimes.

\item The multiplicative \(\log m\) term can be removed from the
\emph{expected null cost} by allowing an unbounded sequential stopping time.
Define the reverse boundary divergence
\begin{equation}\label{eq:formal-results-060}
J_{\rm bnd}(\tau,d)
=
(\tau+d)\log\frac{\tau+d}{\tau}
+(r-d)\log\frac{r-d}{r},
\end{equation}
and
\begin{equation}\label{eq:formal-results-061}
\begin{aligned}
H_0&=\frac{\log(1/\delta)+\log(r/(r-d))}{K_{\rm bnd}(\tau,d)},\\
H_1&=\frac{\log(m/\delta)+\log(1+d/\tau)}{J_{\rm bnd}(\tau,d)}.
\end{aligned}
\end{equation}
There is a sequential hybrid tester with type-I error and every pointwise
type-II error at most \(\delta\), whose expected null declared cost is at most
\begin{equation}\label{eq:formal-results-062}
H_0\sum_{x\in U}\bar c_x
+L C_{\rm ctrl}^{A\setminus U}
+\sum_{i:q_i^\star>0}c_i.
\end{equation}
Under singleton alternative \(x\in U\), its expected boundary-block cost is
at most
\begin{equation}\label{eq:formal-results-063}
H_0\sum_{z\in U\setminus\{x\}}\bar c_z+H_1\bar c_x.
\end{equation}
Thus the boundary block has expected null cost of order
\(d^{-2}\log(1/\delta)\sum_{x\in U}\bar c_x\).
There is no multiplicative \(\log m\); the changed channel still has additive
\(\log m\) delay.  More precisely, if \(Q_0^\star(d)\) denotes the optimal
expected null declared cost in this P19 family and
\(a=\kl(1-\delta,\delta)\), then
\begin{equation}\label{eq:formal-results-064}
\begin{aligned}
\frac{a}{b_2}\sum_{x\in U}\bar c_x
&\le\liminf_{d\downarrow0}d^2Q_0^\star(d)\\
&\le\limsup_{d\downarrow0}d^2Q_0^\star(d)\\
&\le\frac{\log(1/\delta)}{b_2}\sum_{x\in U}\bar c_x.
\end{aligned}
\end{equation}
The remaining multiplicative gap depends only on confidence:
\begin{equation}\label{eq:formal-results-065}
\frac{\log(1/\delta)}{\kl(1-\delta,\delta)}.
\end{equation}
It is independent of topology size and equals approximately
\(1.1305\) at \(\delta=0.05\).  This statement concerns
the exact P19 simple null and singleton alternatives, and does not give a
bounded deterministic horizon.

\item A deterministic declared-cost horizon can be restored without returning
to the fixed-allocation \(m\log m\) penalty.  Fix any
\(\eta\in(0,\delta)\), let \(H_0(\eta)\) and \(H_1(\eta)\) denote the preceding
quantities with \(\delta\) replaced by \(\eta\), and define
\begin{equation}\label{eq:formal-results-066}
B_\eta
=
\frac1{\delta-\eta}
\max_{x\in U}
\left\{
H_0(\eta)\sum_{z\in U\setminus\{x\}}\bar c_z
+H_1(\eta)\bar c_x
\right\}.
\end{equation}
Truncating the sequential boundary block immediately before its declared cost
would exceed \(B_\eta\), and then accepting that block, gives type-I error at
most \(\eta\), pointwise boundary type-II error at most \(\delta\), and the
deterministic total hybrid cost cap
\begin{equation}\label{eq:formal-results-067}
B_\eta
+L C_{\rm ctrl}^{A\setminus U}
+\sum_{i:q_i^\star>0}c_i.
\end{equation}
For unit boundary costs, fixed \(\delta\), and \(\eta=\delta/2\), this cap is
\begin{equation}\label{eq:formal-results-068}
O\!\left(\frac{m+\log m}{d^2}\right)
=O\!\left(\frac{m}{d^2}\right).
\end{equation}
Part~(i) gives a matching \(\Omega(m/d^2)\) lower bound for every deterministic
cap.  Thus adaptive allocation removes the multiplicative \(\log m\) penalty
even with a deterministic horizon, although this Markov truncation is not
claimed to have optimal dependence on \(\delta\) or exact constants.  The
finite-sample certificate can be minimized over the single scalar
\(\eta\in(0,\delta)\); this optimization does not alter the error proof.
\end{enumerate}
\end{theorem}

\begin{proof}
Fix alternative \(x\in A\).  If a queried group contains \(x\) and another
high vertex, the null transcript has a top tape interval of mass \(r\).  Raising
only \(x\)'s new threshold splits mass \(d\) from that interval into a new
\(x\)-specific state, giving \(K_{\rm ctrl}\).  If the group contains \(x\)
but no other high vertex, the existing high--low interval has mass \(\tau\);
the alternative changes masses \((\tau,r)\) to \((\tau+d,r-d)\), giving
\(K_{\rm bnd}\).  Groups omitting \(x\) have zero KL.  Adaptive transcript
change of measure gives one information constraint per \(x\), proving (i).

For (ii), let \(q^\star\) solve the control-cover LP and query group \(S_i\)
\(\lceil Lq_i^\star\rceil\) times.  Within each group, compare old--new patch
deltas only among its high-class members.  Under the null those deltas are
identical pathwise, even though high and low classes differ benignly by
\(\tau\).  Under alternative \(x\), every control-compatible exposure containing
\(x\) becomes nonconstant exactly when
\(U\in(p+s+\tau,p+s+\tau+d]\), an event of probability \(d\).
The cover gives at least \(L\) independent opportunities for each \(x\), so the
miss probability is at most \((1-d)^L\le\delta\).

For (iii), expand the two exact KL expressions.  In the lower LP, scaling
\(z_i=d\,n_i/\kl(1-\delta,\delta)\) makes the coefficient of \(C\) converge to
one and the coefficient of \(R\) converge to zero.  A feasible control-cover
solution gives a uniform upper bound on the scaled optimum.  Because every
\(c_i>0\), any sequence of scaled optimizers therefore has bounded
coordinates; every convergent subsequence has a limit satisfying
\(Cz\ge\mathbf 1\).  Consequently the lower limit of the scaled objective is
at least \(C_{\rm ctrl}\), while the control-cover solution gives the matching
upper limit.  If candidate \(x\) has no control-compatible group, its
information constraint contains only \(K_{\rm bnd}=\Theta(d^2)\), giving the
quadratic lower bound.

For (iv), observe first that a group containing an unmatched candidate
\(x\in U\) cannot contain any other high vertex: otherwise that group would be
a matched control for \(x\).  Hence the LP block for \(x\) is disjoint from
every other high-candidate block and has exact value
\(a\bar c_x/K_{\rm bnd}\).  The Taylor expansion
\begin{equation}\label{eq:formal-results-069}
K_{\rm bnd}(\tau,d)=b_2d^2+b_3d^3+O(d^4)
\end{equation}
therefore contributes
\begin{equation}\label{eq:formal-results-070}
\frac{a\bar c_x}{b_2}d^{-2}
-\frac{ab_3\bar c_x}{b_2^2}d^{-1}+O(1).
\end{equation}
For the covered block, set
\(\epsilon=K_{\rm bnd}/K_{\rm ctrl}=O(d)\) and scale
\(z=K_{\rm ctrl}n/a\).  Its value is
\((a/K_{\rm ctrl})F(\epsilon)\), where \(F(0)\) is precisely
\(C_{\rm ctrl}^{A\setminus U}\).  Positive costs bound every coordinate of
an optimizer uniformly near zero.  Thus the extra term
\(\epsilon Rz\) relaxes each covering constraint by only \(O(\epsilon)\),
and rescaling by \(1+O(\epsilon)\) gives
\(F(\epsilon)=F(0)+O(\epsilon)\).  Since
\(K_{\rm ctrl}=d+O(d^2)\), the covered block contributes
\(aC_{\rm ctrl}^{A\setminus U}d^{-1}+O(1)\).  Summing the disjoint blocks
proves the expansion.

For (v), choose for every \(x\in U\) a cheapest containing group
\(S_{i_x}\).  Because \(x\) has no control-compatible group,
\(S_{i_x}\) contains no other high vertex; since legal groups have size at
least two, it contains a low reference \(y_x\).  On fresh common-uniform
tapes \(R\), form
\begin{equation}\label{eq:formal-results-071}
W_x=\mathbf 1\{P_1(x;R)\ne P_1(y_x;R)\}.
\end{equation}
Under the specified two-class null this is \({\rm Bernoulli}(\tau)\), and under
singleton alternative \(x\) it is \({\rm Bernoulli}(\tau+d)\).  Apply the
P18 threshold test to each \(x\in U\) with error budget \(\delta/m\) and
\(N_U\) rounds.  A union bound controls the boundary block's type-I error by
\(\delta\), while its miss probability under any unmatched alternative is at
most \(\delta/m\).  Independently apply the zero-type-I matched-control
schedule from (ii) to the residual cover on \(A\setminus U\), with pointwise
miss probability at most \(\delta\), and reject if either block rejects.
The declared costs add to the displayed bound.  Groups containing unmatched
candidates cannot contain another high vertex, so they do not create a hidden
cross-coverage term.  The construction relies on the exact equal-high P19 null
and singleton alternatives; it is not valid for an arbitrary heterogeneous
composite null or simultaneous multi-candidate changes.

For (vi), on the cheapest boundary group for each \(x\in U\), let \(L_x(t)\)
be the likelihood ratio of alternative \(x\) against the exact P19 null after
\(t\) fresh rounds.  Process these channels in any fixed order, rejecting if
the active likelihood ratio reaches \(m/\delta\), moving to the next channel
if it reaches \(\delta\), and accepting after every channel has crossed its
lower boundary.  Under the null, Ville's inequality and a union bound give
type-I error at most \(m(\delta/m)=\delta\).  Under alternative \(x\),
\(1/L_x\) is a nonnegative martingale, so the probability of crossing the
lower boundary, and hence the miss probability, is at most \(\delta\).
The channels use disjoint groups and fresh tapes.  Consequently these Ville
bounds remain valid conditionally on the stopping history of every earlier
channel in the fixed processing order; a false upper crossing on an earlier
null channel is already a correct rejection under a singleton alternative and
cannot create a miss.

The log-likelihood increments are bounded.  Their null mean is
\(-K_{\rm bnd}\), and a lower-boundary crossing overshoots by at most
\(\log(r/(r-d))\).  Wald's identity therefore gives
\(\E_0T_x\le H_0\).  Under alternative \(x\), the mean increment is
\(J_{\rm bnd}\), and an upper crossing overshoots by at most
\(\log(1+d/\tau)\), giving \(\E_xT_x\le H_1\).  Summing costs and adjoining
the residual zero-type-I control schedule proves the displayed bounds.  The
exit times are integrable: from every interior likelihood state, a fixed
finite run of identical informative increments exits through one boundary
with probability bounded uniformly away from zero, giving a geometric tail
over such runs.  This also supplies the integrability condition used in
Wald's identity.  Under alternative \(x\), charging every unchanged channel
at its full null stopping cost, including channels that would occur after
\(x\), is an upper bound for every fixed processing order.  The
lower endpoint of the asymptotic bracket is part~(i), because every unmatched
channel contributes \(a\bar c_x/K_{\rm bnd}\); multiplying the lower and upper
bounds by \(d^2\), using \(K_{\rm bnd}=b_2d^2+O(d^3)\), and observing that the
residual control cost is only \(O(d^{-1})\) proves the bracket.

For (vii), run the sequential boundary block from (vi) at internal error
\(\eta\).  Under singleton alternative \(x\), its untruncated declared cost
has expectation at most the expression inside braces defining \(B_\eta\).
Markov's inequality therefore bounds the probability of truncating before the
unbounded test decides by at most \(\delta-\eta\).  Truncation followed by
acceptance cannot increase type-I error, while its pointwise type-II error is
at most \(\eta+(\delta-\eta)=\delta\).  Appending the deterministic residual
control schedule gives the displayed total cap.  Indeed, on the coupled
untruncated path, a pre-query truncation can occur only if the eventual
untruncated declared cost exceeds \(B_\eta\); hence Markov is being applied to
the complete untruncated cost, not to an optional prefix.  With unit costs and
\(\eta=\delta/2\), the numerator defining \(B_\eta\) is
\begin{equation}\label{eq:formal-results-072}
\frac{(m-1)\log(2/\delta)+\log(2m/\delta)}{b_2d^2}
+O(d^{-1}),
\end{equation}
which is \(O(m/d^2)\) at fixed \(\delta\).  Conversely, part~(i) gives null
expected boundary cost at least
\(m\kl(1-\delta,\delta)/K_{\rm bnd}=\Omega(m/d^2)\); every deterministic cost
cap dominates this expectation.
\end{proof}


% Audited October 15 addition: a nuisance-invariant special case and scope fence.
% This is not the missing general P19 composite-null theorem.
\subsection{A nuisance-invariant matched witness and a scope obstruction}
\label{sec:unknown-shared-threshold}

Exact \emph{pathwise} agreement can remove a shared threshold nuisance without
estimating it. Equal \emph{mean} patch effects, however, do not suffice. The
coupling and stopped-KL ingredients are classical
\cite{GlassermanYao1992,NitinawaratAtiaVeeravalli2013}; the proposition records
their consequence for the declared ordered four-output software transcript.
It is a patch-neutrality comparison (or a zero-tolerance edge), not a general
positive-tolerance fairness guarantee or a solution to composite controlled
sensing.

\begin{proposition}[Unknown exactly shared background and a matching obstruction]
\label{prop:unknown-shared-control}
Let one four-way round retain the ordered transcript
\(
X=(P_0(x;U),P_1(x;U),P_0(y;U),P_1(y;U))
\), where \(P_j(z;U)=\mathbf 1\{U\le p_{j,z}\}\), the same
\(U\sim\mathrm{Unif}[0,1]\) is used for all four outputs, and tapes are
independent between rounds. Fix \(\delta\in(0,1/2)\) and known \(d>0\).
\begin{enumerate}[label=(\roman*)]
\item Let
\(
\mathcal N_d=\{(p,h):0<p<h<h+d<1\}
\).
For \((p,h)\in\mathcal N_d\), the null thresholds are
\((p,p)\to(h,h)\), whereas the corresponding alternative is
\((p,p)\to(h+d,h)\). The patch-delta discrepancy witness has probability
zero under every null and exactly \(d\) under every corresponding alternative.
Consequently,
\begin{equation}\label{eq:unknown-shared-horizon}
L(d,\delta)=\left\lceil\frac{\log\delta}{\log(1-d)}\right\rceil
\end{equation}
independent rounds have uniform type-I error zero and type-II error at most
\(\delta\), at a cost of \(4L(d,\delta)\) concrete program executions.

Conversely, consider a sequential tester adapted to the independent-round
transcript filtration, terminating with a decision, whose two errors are at
most \(\delta\) uniformly over the paired null/alternative families. For every
fixed \((p,h)\in\mathcal N_d\), if its null expected stopping time is finite,
then
\begin{equation}\label{eq:unknown-shared-lower}
\E_{0,p,h}[T]\ge
\frac{\kl(1-\delta,\delta)}
{(1-h)\log\bigl((1-h)/(1-h-d)\bigr)}.
\end{equation}
Thus, for fixed \(\delta\) and nuisance restricted to any fixed compact subset
of \(\{(p,h):0<p<h<1\}\), the minimax null-expected round complexity is
\(\Theta(d^{-1})\) as \(d\downarrow0\). This is only a null-side converse;
the reverse single-round KL is infinite.

\item Consider instead the fair null with old thresholds \((p,p+\eta)\) and
new thresholds \((p+s,p+s+\eta)\), where
\begin{equation}\label{eq:mean-matching-domain}
0<p<p+\eta<p+s<p+s+\eta<1.
\end{equation}
Both vertices have expected patch effect \(s\), and both versions have edge
disparity \(\eta\le\tau\). Nevertheless, their tape-wise patch deltas disagree
with probability exactly \(2\eta\) per round. Applying the zero-witness
procedure for \(L(d,\delta)\) rounds to this fair null has false-rejection
probability
\begin{equation}\label{eq:mean-matching-failure}
1-(1-2\eta)^{L(d,\delta)},
\end{equation}
which tends to one for fixed \(\eta>0\) as \(d\downarrow0\). Equality of mean
patch effects therefore cannot replace pathwise matching.
\end{enumerate}
\end{proposition}

\begin{proof}
For part (i), in the coordinate order of the proposition, the complete
single-round laws are
\begin{align}
P_{0,p,h}:&\quad
1111:p,\qquad 0101:h-p,\qquad 0000:1-h,
\label{eq:unknown-shared-null-atoms}\\
P_{1,p,h}:&\quad
1111:p,\qquad 0101:h-p,\qquad 0100:d,\qquad
0000:1-h-d.
\label{eq:unknown-shared-alt-atoms}
\end{align}
The new atom \(0100\) is exactly the patch-delta discrepancy event. This gives
the two witness probabilities and the miss probability \((1-d)^L\).
Moreover, \(\operatorname{supp}(P_{0,p,h})\) is a strict subset of
\(\operatorname{supp}(P_{1,p,h})\), and hence
\begin{equation}\label{eq:unknown-shared-kl-directions}
D(P_{0,p,h}\Vert P_{1,p,h})
=(1-h)\log\!\frac{1-h}{1-h-d},
\qquad
D(P_{1,p,h}\Vert P_{0,p,h})=\infty.
\end{equation}
Under the null, the one-round log-likelihood increment is bounded. The stopped
transcript chain rule and data processing through the terminal decision give
\(\E_{0,p,h}[T]D(P_{0,p,h}\Vert P_{1,p,h})\ge
\kl(1-\delta,\delta)\), proving \Cref{eq:unknown-shared-lower}.
Uniformly when \(1-h\) is bounded away from zero,
\(
D(P_{0,p,h}\Vert P_{1,p,h})=d+d^2/(2(1-h))+O(d^3)
\), while \(L(d,\delta)=\Theta(d^{-1})\) at fixed \(\delta\).

For part (ii), the two patch-delta indicators correspond to intervals
\((p,p+s]\) and \((p+\eta,p+s+\eta]\). Under
\Cref{eq:mean-matching-domain}, their symmetric difference consists of two
disjoint intervals of length \(\eta\), so its measure is \(2\eta\).
Independence between rounds proves \Cref{eq:mean-matching-failure} and its
limit. The endpoint mean shifts are both \(s\), while both the old and new
inter-endpoint disparities equal \(\eta\).
\end{proof}


\begin{proposition}[A fixed-allocation familywise penalty]
\label{prop:boundary-familywise-fixed}
Consider \(m\) disjoint, unit-cost P19 boundary channels, and suppose a
nonadaptive design takes deterministic query counts \(n_1,\ldots,n_m\).
If a test has type-I error and every pointwise type-II error at most
\(\delta\), write \(n_{(1)}\le\cdots\le n_{(m)}\).  Then, for every
\(k\in\{1,\ldots,m\}\),
\begin{equation}\label{eq:formal-results-073}
\begin{aligned}
n_{(k)}&\ge\frac{\log(1+D_\delta k)}{\rho_2},\\
D_\delta&=\frac{(1-2\delta)^2}{\delta(1-\delta)},\\
\rho_2&=\log\!\left(1+d^2(\tau^{-1}+r^{-1})\right).
\end{aligned}
\end{equation}
Consequently,
\begin{equation}\label{eq:formal-results-074}
\sum_{x=1}^m n_x
\ge
\frac1{\rho_2}\sum_{k=1}^m\log(1+D_\delta k)
=\Omega\!\left(\frac{m\log m}{d^2}\right)
\end{equation}
for fixed \(\delta,\tau,r\).  Hence the \(\log m\) factor in the fixed
boundary allocation of part~(v) is not merely a union-bound artifact.  The
proposition does not cover adaptive allocations; part~(vi) exploits precisely
that distinction for expected null cost.
\end{proposition}

\begin{proof}
For one boundary round, the null masses \((\tau,r)\) change to
\((\tau+d,r-d)\), so its chi-square divergence is
\(d^2(1/\tau+1/r)\).  After \(n_x\) independent rounds the divergence is
\(\exp(\rho_2n_x)-1\).  Fix the \(k\) least-sampled channels and mix their
singleton alternatives uniformly.  Disjointness makes the mixture
chi-square divergence from the null
\begin{equation}\label{eq:formal-results-075}
\frac1{k^2}\sum_{j=1}^k
\bigl(\exp(\rho_2n_{(j)})-1\bigr)
\le
\frac{\exp(\rho_2n_{(k)})-1}{k}.
\end{equation}
Data processing through the terminal decision lower-bounds the same divergence
by the Bernoulli decision divergence
\(D_\delta\), because rejection has probability at most \(\delta\) under the
null and at least \(1-\delta\) under the mixture.  Rearranging gives the
rankwise bound; summing over \(k\) proves the claim.
\end{proof}

